%% Sometimes when a section can't be nicely modelled with the \entry[]... mechanism; hack our own
\makerubrichead{Research Experiences}

%% Assuming you've already given \addbibresource{own-bib.bib} in the main doc. Right? Right???
\nocite{*}

%% If you just want everything in one list
% \printbibliography[heading={none}]

% \printbibliography[heading={subbibliography},title={Journal Articles},type=article]

% \printbibliography[heading={subbibliography},title={Conference Proceedings},type=inproceedings]

% \printbibliography[heading={subbibliography},title={Books and Chapters},filter={booksandchapters}]


\noindent\textbf{Continuous Planning \& Real-time Adaptation for Collaborative LTL Tasks} \hfill \textit{Jul. 2025 -- Present}
\begin{itemize}[parsep=0.7ex, leftmargin=1.5em]
  \item Built a multi‑heterogeneous robot simulation platform (Unitree Go2, Scout) in Gazebo; validated algorithm robustness via kinematics comparisons.
  \item Designed a real‑time adaptive strategy for LTL task decomposition; currently transferring to physical robot experiments.
  % \item \textbf{Expected output:} 1 journal paper in preparation for \textit{IJRR}.
\end{itemize}

\noindent\textbf{LLM‑based Multi‑robot Cooperative Task Planning in Open Worlds} \hfill \textit{Jan. 2026 -- Present}
\begin{itemize}[parsep=0.7ex, leftmargin=1.5em]
  \item Completed literature review and offline planning code; currently integrating AI2‑THOR simulation to test task allocation.
  \item Developing an online replanning module using LLM reasoning to handle environmental changes.
  % \item \textbf{Expected output:} 1 conference paper in preparation for \textit{ICRA}.
\end{itemize}